\chapter{The MDL Growth Law}\label{ch:mdl}

\begin{workplan}{\Recast}
\wpgoal{Keep Construction~7 and give it its V2 placement: this is the
\emph{maintenance cost} of the sleep phase, and it is the only place cost
lives in the entire architecture.}

\wpsource{This chapter's existing text;
\texttt{MATH\_TOOLBOX.md}~III.4 (Construction~7, the per-direction code, the
measurement); \texttt{DOCS/DERIVATION\_MDL\_GROWTH\_LAW.md};
\texttt{TATIANA\_V1\_IDENTITY\_MAP.md}~\S6 (placement, and the one free
parameter).}

\wpsteps{
  \item \textbf{Keep the superseded law as a worked failure.} The original had
    a numerator overcounting by a factor of $k$ (fixed by holonomy
    propagation --- only the first leg needs describing, the rest are
    determined) and, fatally, a denominator that did not depend on concept
    size at all, making a zero-dimensional concept always optimal: ``attach
    always, with a concept that carries nothing.'' Keep this; it shows exactly
    how a plausible law can invert its own purpose.
  \item \textbf{State Construction~7's repair.} A per-direction code: fix a
    direction in the limit, take its $k$ readings around a traversal, and ask
    whether they agree. The saving per traversal follows from the intraclass
    correlation, and the law becomes a threshold on recurrence count that is
    never met when the saving is non-positive.
  \item \textbf{Report the measurements.} Within $0.042$ bits/traversal of
    exact quantised-Gaussian entropy; over-prices its own reference code so
    the threshold is \emph{conservative} --- it under-attaches, never
    over-attaches; and a shuffled-readings control drops the saving from
    $+5.289$ to $-0.437$ bits, so the saving measures the recurrence and not
    the coder. \Measured.
  \item \textbf{Keep the restored consequences.} Size-dependence returns
    per-direction; the subspace-selection rule (rank by reliability, keep a
    prefix) is now a \emph{theorem} rather than a proposal; and a reliability
    floor exists independent of frequency.
  \item \textbf{Keep the open self-reference problem.} Taking bits-per-real as
    $\tfrac12\log n$ is self-referential but resolvable by iteration ---
    \emph{except} that it degenerates exactly where the growth law lives: at
    $n = 1$ the model is free and over-generation returns. Do not adopt it as
    derived-therefore-safe. State the alternatives.
  \item \textbf{Keep the stated narrowing.} The law prices what was there, not
    which succession fired --- a strict under-estimate, making it conservative
    by construction rather than by luck.
  \item \textbf{Add the V2 placement paragraph.} This cost is consulted during
    \emph{sleep} only. It never bids against growth inside a single tick.
    Name the exchange rate between deformation cost and new-cell cost as the
    one genuine free parameter of the design, and record that acquiring a
    second one is evidence the two-phase separation has failed.
}

\wpdone{The law is stated with its measurements and its open self-reference
problem, and its placement in the sleep phase is explicit.}
\end{workplan}

The question this chapter answers: \emph{when is a new concept worth
inventing?} Earlier designs answered it with a threshold someone chose.
This one derives a threshold instead.

\section{Two-part description length}

\begin{intu}\Intu
A memory is a compression of experience. Minimum description length
makes that literal: the total cost of an explanation is what it takes to
write down the model plus what it takes to write down the data
\emph{given} the model. Adding a concept always makes the first term
worse. It is worth doing exactly when it makes the second term better by
more.
\end{intu}

\begin{definition}[Two-part MDL]\label{def:mdl}
\[
  L_{\text{total}}
  =\underbrace{L(\text{model})}_{\text{the memory}}
  +\underbrace{L(\text{data}\mid\text{model})}_{\text{the experience}} .
\]
The \emph{model} is the complex $\CC$ together with the sheaf $\Fs$ ---
what the system \emph{is}. The \emph{data} is the observed record: the
sequence of assemblies and the succession relation read off it --- what
the system has \emph{seen}.
\end{definition}

Neither term alone is the criterion. Minimising the model term alone
forbids all growth; minimising the data term alone is achieved by
memorising everything. The criterion is the sum.

\section{The model term}

\begin{proposition}[Cost of coning a cycle]\label{prop:modelterm}
\Proved\ Let $\gamma$ be a $1$-cycle of length $k$, and cone it with a
new vertex $w$ carrying a stalk of dimension $d_w$. Writing $b$ for bits
per real number, the added model cost is
\begin{equation}\label{eq:modelterm}
  \Delta L(\text{model})
  =b\,d_w\Bigl(1+\sum_{e\in\gamma}d_e\Bigr).
\end{equation}
\end{proposition}
\begin{proof}
Three kinds of new object are described. The vertex with its stalk costs
$d_w\cdot b$ bits. Each of the $k$ restriction maps out of $\stalk{w}$
is a $d_e\times d_w$ matrix, costing $d_w d_e b$ bits, summed over
$e\in\gamma$. The $k$ new triangles cost \emph{nothing}: they are
determined by the cone construction, not described separately. Summing
gives \eqref{eq:modelterm}.
\end{proof}

\begin{remark}[Two consequences, both load-bearing]
First, the cost is \emph{linear in $d_w$}. The dimension of the new
stalk is the dominant free quantity, so whatever fixes $d_w$ fixes most
of the cost --- which is what Section~\ref{sec:conesize} is for.
Second, \emph{triangles are free}. They are implied by the cone, so the
$k$ new $2$-cells cost nothing to describe. This is a real asymmetry,
not a bookkeeping convenience, and it is why coning is cheap relative to
what it buys.
\end{remark}

\section{The data term and the law}

\begin{proposition}[Data-term saving]\label{prop:dataterm}
\Proved\ Let the cycle $\gamma$ occur $n$ times in the record. Without
the concept, each traversal is coded as $k$ separate transitions, each
costing its full surprisal; with the concept it is a single named event.
Then
\begin{equation}\label{eq:dataterm}
  \Delta L(\text{data}\mid\text{model})=-n\bigl(c_{\text{old}}-c_{\text{new}}\bigr),
\end{equation}
where $c_{\text{old}}$ is the cost of coding one traversal the long way
and $c_{\text{new}}$ the cost of naming it.
\end{proposition}

\begin{theorem}[The growth law]\label{thm:growthlaw}
\Proved\ Attach the concept if and only if $\Delta L_{\text{total}}<0$,
that is
\begin{equation}\label{eq:growthlaw}
  \boxed{\ n\;>\;
  \frac{b\,d_w\bigl(1+\sum_{e\in\gamma}d_e\bigr)}
       {c_{\text{old}}-c_{\text{new}}}\ }
\end{equation}
\end{theorem}
\begin{proof}
Add \eqref{eq:modelterm} and \eqref{eq:dataterm} and require the sum
negative; rearrange, using $c_{\text{old}}>c_{\text{new}}$.
\end{proof}

\begin{intu}\Intu
Read what this is: a \emph{recurrence threshold that is computed, not
chosen}. It is the ratio of what a concept costs to store against what
it saves per use --- the number of times a pattern must repeat before it
is worth naming.
\end{intu}

Three consequences, two of which were not asked for:

\begin{enumerate}
  \item $n$ is a number the engine can print, rather than a parameter a
        human sets.
  \item \textbf{It orders the holes.} When several harmonic classes
        compete, fill them in decreasing $\abs{\Delta L_{\text{total}}}$.
        This removes a decision that had been silently delegated to a
        human.
  \item \textbf{It predicts non-attachment.} A one-off inconsistency is
        \emph{noise to be tolerated}, not a concept to be invented. The
        earlier threshold-based law had no way to express this, which is
        exactly how it would have over-generated.
\end{enumerate}

\chapter{The Coned Concept}\label{ch:cone}

\begin{workplan}{\Keep}
\wpgoal{Keep Construction~6 essentially as written --- it is the
most-verified construction in the record --- and add the one upgrade V2 needs:
its output dimension is what moves the invariant, so this chapter and
Chapter~\ref{ch:identity} must agree exactly.}

\wpsource{This chapter's existing text;
\texttt{MATH\_TOOLBOX.md}~III.3 (derivation, two self-corrections, the
measurements, provenance); \texttt{TATIANA-V2/MATH/phase0\_invariant\_dimension.md}
(the dimension confirmed two independent ways).}

\wpsteps{
  \item \textbf{Keep the derivation and its punchline.} Consistency on the new
    edges forces the value at $w$ to \emph{generate} the section over the whole
    cycle, and substituting into the old edges gives verbatim the categorical
    cone condition. The cone condition was not imposed --- it fell out of
    demanding the topological cone do its job. \Proved.
  \item \textbf{Keep the identification.} $\Fs(w) := \lim D_\gamma \cong
    \Hzero(\gamma;\Fs|_\gamma)$: the new concept is \emph{not} a summary of the
    cycle's members but their \emph{compatible part}. This sentence is the
    chapter.
  \item \textbf{Keep both self-corrections visible.} First, an earlier draft
    argued the limit is ``the smallest object admitting the required maps'' ---
    backwards; every cone factors uniquely through the limit, so it is the
    \emph{largest} non-redundant choice, a ceiling and not a floor. Second, the
    ``refusal branch'' (nothing consistent, so the construction declines) is
    logically real but \emph{structurally unreachable}: a cycle carries
    harmonic mass if and only if it has sections.
  \item \textbf{Keep the upgrade from preference to necessity.} The coned
    complex is a cochain complex \emph{if and only if} $\Fs(w)$ is a cone over
    the diagram. Construction~6 is therefore not the best choice among several
    --- it is the only choice for which the coned object is a sheaf at all. The
    universal property was never a preference; it was the existence condition.
    \Proved, and measured: an arbitrary $\Fs(w)$ gives $1.756$, the limit
    gives $4.4 \times 10^{-16}$.
  \item \textbf{Keep the parity observation.} Reachability of the degenerate
    case is governed by the parity of $d$, since every element of $SO(\text{odd})$
    fixes an axis. This is a fact about this architecture specifically, because
    its restriction maps are orthogonal by construction.
  \item \textbf{Keep the provenance paragraph, including the direction.} The
    universal-property route to a new concept is Goguen's, not this project's
    --- and the variance is \emph{opposite}: his blends are colimits, a sheaf's
    restriction maps force a limit. ``Copying him would have gotten it
    backwards.'' What is genuinely original is that the sheaf's own restriction
    maps supply the diagram, so nothing new is introduced beyond what the
    memory already held.
  \item \textbf{Add the invariant link.} $\dim\Fs(w) = \dim\Hzero(\gamma;
    \Fs|_\gamma)$ is exactly the amount by which $\mathrm{K}(\mathbb{K})$
    changes, known \emph{before} the growth decision, from the same
    computation that produced the address. Cross-reference
    Chapter~\ref{ch:identity} and report that both routes to this number have
    been computed and agree. \Measured.
  \item \textbf{Record what is still not derivable.} The label for the new
    concept is not derivable from a diagram. Name this as the one place
    genuine novelty enters rather than being computed, and connect it to
    Chapter~\ref{ch:encoder}.
}

\wpdone{The construction is proved to be forced rather than chosen, and its
output dimension is stated identically here and in Chapter~\ref{ch:identity}.}
\end{workplan}

This is the central construction of the later work: the stalk of a
newly-grown concept, \emph{derived} rather than chosen.

\section{The topological cone}

\begin{construction}[Coning a cycle]\label{con:conetopo}
Let $\gamma=(v_1,e_1,v_2,\dots,v_k,e_k,v_1)$ be a $1$-cycle carrying
harmonic mass, with $e_i=\{v_i,v_{i+1}\}$ (indices mod $k$). Attach a
vertex $w$, edges $f_i=\{w,v_i\}$, and triangles
$t_i=\{w,v_i,v_{i+1}\}$.
\end{construction}

\begin{proposition}[The cone kills the cycle]\label{prop:conekills}
\Proved\ With $\partial t_i=e_i-f_{i+1}+f_i$,
\[
  \partial\Bigl(\sum_i t_i\Bigr)
  =\sum_i e_i+\sum_i(f_i-f_{i+1})=\gamma,
\]
since the $f$ terms telescope to zero. Hence $\gamma$ becomes a
boundary.
\end{proposition}

\section{The stalks, chosen minimally}

\begin{construction}[New edge stalks]\label{con:conestalks}
Set
\[
  \stalk{f_i}:=\stalk{v_i},
  \qquad
  \rest{v_i\triang f_i}:=\Id .
\]
\end{construction}

\begin{hon}\Hon
This is an \emph{engineering choice}, and a deliberately empty one. Any
distortion introduced here would be a free parameter, and whatever the
construction then achieved could be attributed to a cleverly chosen edge
stalk rather than to the concept. All the content is forced into
$\stalk{w}$ and the maps out of it, which is where it can be reasoned
about.
\end{hon}

The only remaining unknowns are $\stalk{w}$ and the legs
$p_i:=\rest{w\triang f_i}:\stalk{w}\to\stalk{v_i}$.

\section{The cone condition is forced, not imposed}

\begin{proposition}[Derivation of the cone condition]\label{prop:conecond}
\Proved\ Write $r_i^-:=\rest{v_i\triang e_i}$ and
$r_i^+:=\rest{v_{i+1}\triang e_i}$. For a $0$-cochain $x$:
\begin{enumerate}[label=(\roman*)]
  \item $\cob x$ vanishes on the new edges if and only if
        $x_{v_i}=p_i(x_w)$ for every $i$;
  \item given (i), $\cob x$ vanishes on the old edges for
        \emph{every} $x_w\in\stalk{w}$ if and only if
        \begin{equation}\label{eq:conecond}
          \boxed{\ r_i^+\circ p_{i+1}=r_i^-\circ p_i
          \quad\text{for all }i\ }
        \end{equation}
\end{enumerate}
\end{proposition}
\begin{proof}
(i) On the new edge $f_i$, using
Construction~\ref{con:conestalks},
\[
  (\cob x)_{f_i}
  =\rest{w\triang f_i}x_w-\rest{v_i\triang f_i}x_{v_i}
  =p_i(x_w)-x_{v_i},
\]
which vanishes exactly when $x_{v_i}=p_i(x_w)$.

(ii) Substituting $x_{v_i}=p_i(x_w)$ into the old edge $e_i$,
\[
  (\cob x)_{e_i}=r_i^+p_{i+1}(x_w)-r_i^-p_i(x_w)
  =\bigl(r_i^+p_{i+1}-r_i^-p_i\bigr)(x_w).
\]
A linear map vanishes on all of $\stalk{w}$ if and only if it is the
zero map, which is \eqref{eq:conecond}.
\end{proof}

\begin{intu}\Intu
Part (i) says something worth pausing on: \textbf{the value at $w$
generates the section over the whole cycle}. That is what it means for a
concept to explain a loop --- the loop stops being $k$ independent facts
and becomes one.
\end{intu}

Condition \eqref{eq:conecond} is verbatim the definition of a
\emph{cone} over the diagram
\[
  D_\gamma:\quad
  \stalk{v_1}\xrightarrow{r_1^-}\stalk{e_1}\xleftarrow{r_1^+}
  \stalk{v_2}\xrightarrow{r_2^-}\stalk{e_2}\xleftarrow{\ \ }\cdots
\]
with apex $\stalk{w}$ and legs $p_i$. The categorical cone condition was
not imposed on the construction; it fell out of demanding that the
topological cone do its job.

\section{The construction}

\begin{construction}[The coned concept]\label{con:coned}
\[
  \boxed{\ \stalk{w}:=\varprojlim D_\gamma,
  \qquad p_i:=\text{the limit's projections.}\ }
\]
\end{construction}

\begin{proposition}[Concrete description]\label{prop:conelim}
\Proved\ For finite-dimensional stalks,
\[
  \varprojlim D_\gamma
  =\Bigl\{(x_1,\dots,x_k)\in\bigoplus_i\stalk{v_i}\ \Big|\
    r_i^-x_i=r_i^+x_{i+1}\ \forall i\Bigr\}
  \;\cong\;H^0\bigl(\gamma;\Fs|_\gamma\bigr).
\]
\end{proposition}
\begin{proof}
The limit of a finite diagram of vector spaces is the subspace of the
product on which all the diagram's maps agree, which is the displayed
equaliser. Comparing with Definition~\ref{def:coboundary} restricted to
the subcomplex $\gamma$, the agreement conditions are exactly
$\cob x=0$, so the space is $\ker\cob|_\gamma=H^0(\gamma;\Fs|_\gamma)$.
\end{proof}

\begin{intu}\Intu
In words: the new concept \emph{is} everything that was consistent
around the loop but could not be glued. It does not summarise the
cycle's members; it \emph{is} the compatible part of them.
\end{intu}

Three properties, each of which had been written down separately as a
requirement:

\begin{center}
\begin{tabular}{@{}p{3.4cm}p{8.4cm}@{}}
\toprule
\textbf{Property} & \textbf{Why it matters}\\
\midrule
canonical --- unique up to unique isomorphism & no hand-made choice
enters the construction\\
path-dependent --- depends on $D_\gamma$, i.e. on which cycle history
produced & precisely the property whose absence caused an earlier
formal-concept-analysis approach to be rejected\\
degenerate case & if $H^0(\gamma;\Fs|_\gamma)=0$ the limit is the zero
space --- but see Proposition~\ref{prop:unreachable}\\
\bottomrule
\end{tabular}
\end{center}

\section{The refusal branch is unreachable}

\begin{proposition}[Harmonic mass implies sections]\label{prop:unreachable}
\Proved\ On a $k$-cycle with vertex and edge stalks all of dimension
$n$, and with no $2$-cells,
\[
  \dim H^1(\gamma;\Fs|_\gamma)=\dim H^0(\gamma;\Fs|_\gamma).
\]
Consequently the cycle carries harmonic mass if and only if it has
nonzero sections.
\end{proposition}
\begin{proof}
For a $1$-dimensional complex the Euler characteristic satisfies
$\chi=\dim\Czero-\dim\Cone=\dim H^0-\dim H^1$. On a $k$-cycle with all
stalks of dimension $n$ there are $k$ vertices and $k$ edges, so
$\dim\Czero=\dim\Cone=kn$ and $\chi=0$, giving the equality. With no
$2$-cells the harmonic space \emph{is} $H^1$.
\end{proof}

\begin{hon}\Hon
An earlier draft claimed the degenerate case is meaningful --- ``nothing
was consistent around that loop, so the construction refuses rather than
inventing''. True, and it never happens. By
Proposition~\ref{prop:unreachable}, $H^0=0$ implies no harmonic mass,
which implies no growth address, which implies the growth law never
fires on that cycle at all. Confirmed numerically: the degenerate case
reports that the harmonic space is already empty before coning, so there
is nothing to kill.

In general the gap is $\dim H^1-\dim H^0=k(n_e-n_v)$, so the branch
\emph{could} become reachable if edge stalks are ever given a different
dimension from vertex stalks. They are not, currently.
\end{hon}

\section{The limit is a ceiling, not a floor}\label{sec:conesize}

\begin{hon}\Hon
\Refuted\ An earlier draft justified
Construction~\ref{con:coned} by asserting that the limit is ``the
smallest object admitting the required maps''. That is backwards, and
the error is worth keeping visible.
\end{hon}

\begin{proposition}[Universality bounds from above]\label{prop:ceiling}
\Proved\ Every cone over $D_\gamma$ factors uniquely through the limit.
If a candidate apex $\stalk{w}'$ has jointly injective legs --- i.e.
carries no data that no restriction map ever sees --- then that
factoring map is injective, so
\[
  \dim\stalk{w}'\ \le\ \dim\varprojlim D_\gamma .
\]
\end{proposition}
\begin{proof}
The universal property of the limit gives a unique
$\phi:\stalk{w}'\to\varprojlim D_\gamma$ commuting with the legs. If
$\phi(z)=0$ then every leg $p_i'(z)=0$; joint injectivity of the legs
then forces $z=0$, so $\phi$ is injective and the dimension inequality
follows.
\end{proof}

So the limit is the \emph{largest non-redundant} choice, not the
smallest --- a ceiling on what a new concept can usefully carry. The two
determinations are therefore genuinely separate, and both are derived:

\begin{enumerate}
  \item \textbf{Consistency fixes the shape.} $\stalk{w}$ must be a cone
        over $D_\gamma$; the limit is the universal one.
  \item \textbf{MDL fixes the size.} By \eqref{eq:modelterm} the cost is
        linear in $d_w$, so the optimum is a \emph{subspace} of the
        limit: keep the directions that pay for themselves, discard the
        rest.
\end{enumerate}

\chapter{Retrieval: Rank, Do Not Threshold}\label{ch:retrieval}

\begin{workplan}{\Recast}
\wpgoal{Keep the empirical result, which is large and well-measured, but
change what the chapter is \emph{for}: V2's own diagnosis is that retrieval
cannot live in a metric at all, so this chapter becomes the evidence for that
conclusion rather than a description of a component.}

\wpsource{This chapter's existing text;
\texttt{MATH\_TOOLBOX.md}~V.1--V.2 (the bug, the fix, the self-audit), VII.1
(the six measurements that together say retrieval needs a generative process).}

\wpsteps{
  \item \textbf{Keep the bug and its magnitude.} A relevance threshold
    compared against squared Bures--Wasserstein distance, with an unbounded
    scan and no limit, admitted $0.31\%$ of ground-truth dependencies; at the
    shipped default, literally zero of three thousand ticks retrieved a pair.
    \Measured.
  \item \textbf{Keep all three root causes}, since each is instructive: it was
    a near-duplicate filter rather than a relevance filter; embeddings sit on a
    large common-mode pedestal that any absolute threshold mostly measures; and
    the epistemic term is dimension-extensive while the semantic term is
    dimension-intensive, so summing them lets the former dominate by two orders
    of magnitude regardless of meaning.
  \item \textbf{Keep the fix and the honest framing of its size.} Rank, do not
    threshold; rank on cosine, not squared distance (stored means are not
    unit-norm, so squared-distance ranking silently penalises general
    concepts). Worth $0.31\% \to 46.6\%$. State the lesson explicitly: the
    whole design argument was about a four-point effect while the actual fix
    was worth $140\times$.
  \item \textbf{Keep the self-audit.} Two claims about the fix were themselves
    overstated and walked back --- an optimum read off a flat sweep, and a
    parameter called ``derived'' that was an unmeasured budget coincidence.
    ``Two independent constraints landing on the same number is exactly the
    tidiness motivated reasoning produces.'' Keep this sentence.
  \item \textbf{Keep the refuted alternative.} Diffusion/heat-kernel retrieval,
    twelve configurations, none beating raw cosine, degrading monotonically
    with diffusion time --- the pedestal makes the walk near-uniform so it
    amplifies hubness rather than removing it. \Refuted.
  \item \textbf{Add the V2 recast.} State the conclusion this chapter is now
    evidence for: retrieval cannot live in the metric alone, cannot live in a
    single relation (citation and embedding similarity agree under half the
    time), and cannot live in the corpus. What is missing is a process that
    decides what to retrieve \emph{about} --- which in V2 is the wake phase's
    drive, not a lookup. Reframe accordingly and forward-reference
    Chapter~\ref{ch:twophase}.
  \item \textbf{Carry the standing warning.} A simulator produces the topology
    its rules imply. Moving from ``assume a structure and test for it'' to
    ``generate the structure'' is a stronger design and \emph{also} a way to
    hide assumptions inside an update rule. No claim about emergent structure
    without a validation battery built \emph{before} the generator.
}

\wpdone{The chapter reads as the argument that retrieval needs a process
rather than a metric, with the measurements as evidence rather than as the
point.}
\end{workplan}

\begin{intu}\Intu
Retrieval decides which stored concepts an organ sees. A threshold on
distance answers ``which concepts are close enough?'', which sounds
right and is not: it makes the number of results depend on how the
embedding space happens to be scaled, and it silently returns nothing at
all when the query lands in a sparse region.
\end{intu}

\begin{construction}[Bounded rank query]\label{con:topk}
Replace the $\varepsilon$-ball query with a bounded rank query returning
the top $k$ stored concepts, ranked by \emph{cosine similarity} to the
query mean, using a bounded max-heap of size $k$. Zero-norm vectors are
dropped rather than scored, since a zero vector has no direction.
\end{construction}

\begin{proposition}[Why cosine and not squared distance]\label{prop:cosine}
\Proved\ Expanding,
\[
  \norm{\mu_q-\mu_i}^2
  =\norm{\mu_q}^2-2\inner{\mu_q}{\mu_i}+\norm{\mu_i}^2 .
\]
For a fixed query, $\norm{\mu_q}^2$ is constant and drops out of any
ranking. The term $\norm{\mu_i}^2$ does \emph{not}. Hence squared
distance reproduces the cosine ranking if and only if every stored mean
is unit-norm.
\end{proposition}

\begin{hon}\Hon
Stored means are \emph{not} unit-norm. By
Construction~\ref{con:piv}, $\pi_v$ is a weighted centroid of unit
vectors, so its norm lies below $1$ and falls further the more spread
the concepts it summarises. Ranking by squared distance would therefore
apply a per-concept penalty \emph{proportional to how broad a concept
is}, pushing exactly the general concepts down the list for a reason
unrelated to the query.

There is also an evidential argument, which is decisive on its own: the
measured recall figure against which the retrieval fix is accepted was
obtained on a \emph{cosine} ranking. Ranking by squared distance is a
different ordering carrying an unmeasured bias, which would have made
the acceptance criterion meaningless.
\end{hon}

\begin{remark}[Cost]
Scoring needs only $\mu$; the low-rank factor $U$, the floor $D$ and the
stored reasoning text are not needed unless the row enters the heap.
Reading $\mu$, scoring, and deserialising the rest only on entry drops
the per-row cost from $O(dk^2+k^3)$ to $O(d)$ for the large majority of
rows.
\end{remark}

\begin{hon}\Hon
$k$ is an \emph{engineering choice} ($k=24$), not a derived quantity.
What changed is that it is now constrained on both sides --- too small
starves the organ, too large defeats the point of retrieval --- rather
than being an unexamined constant. Saying so is the difference between a
parameter and a magic number.
\end{hon}
